\documentclass[a4paper,11pt]{article}
\usepackage[scale={0.8,0.9},centering,includeheadfoot]{geometry}
\usepackage{amstext}
\usepackage{listings}
\begin{document}

\section*{PC prior for the range in PRW2}

This is the PC prior for the (log-)range for the PRW2 latent model,
and is only ment for that model.

\subsection*{Parametrization}

This is the prior for the range for the PRW2 model
\begin{displaymath}
    x_t = 2\gamma x_{t-1} - \gamma^{2} x_{t-2} + \epsilon_t, \qquad
    |\gamma| < 1
\end{displaymath}
where $\epsilon_t$ is iid zero mean Gaussian noise. This process is
made stationary with unit variance, and the range, the 'distance' or 'time' to
small correlation, is set to be
\begin{displaymath}
    \rho = \exp\left( - \frac{\sqrt{12} h}{r} \right)
\end{displaymath}
where $h$ is (typical) distance or time, `step-size', between the knots, here $h=1$.
The correlation function of the process is found to be
\begin{displaymath}
    \rho(h) = \left(1  + h \frac{1-\gamma^{2}}{1+\gamma^{2}}\right)\gamma^{h},
    \qquad h=0, 1, 2, \ldots
\end{displaymath}
The PC prior is derived with the limiting case, $\gamma\rightarrow
1^{-}$, as the base model. The distance if found to be
\begin{displaymath}
    d(\gamma) = \sqrt{2\text{KLD}} = \sqrt{3 - 4\rho(1) + \rho(2)} + o(1)
\end{displaymath}
In distance scale, the PC-prior is exponential with rate $\lambda$,
\begin{displaymath}
    \pi(d) = \lambda \exp(-\lambda d), \qquad \lambda > 0.
\end{displaymath}
The calibration of $\lambda$, is done so that 
\begin{displaymath}
    \text{Prob}(r < r_0) = \alpha, \quad r_0 > 0, \; 0 < \alpha < 1
\end{displaymath}
for given $r_0$ and $\alpha$. This prior is also applied to PRW2 for
irregular locations derived from the corresponding limiting process in
continous time.

\subsection*{Specification}

This prior is specified as follows, as
\begin{verbatim}
    hyper = list(range = list(prior = "pc.prw2.range",
                              param = c(r0, alpha, h, lambda)))}
\end{verbatim}
where \texttt{r0}$=r_0$, \texttt{alpha}$=\alpha$, \texttt{h}$=h$ (the
step-size), \texttt{lambda}$=\lambda$ (the step-size).

The default value used for $r_0$, if initially $r_0 \le 0$, is
\begin{verbatim}
    r0 <- diff(range(loc)) / 4
\end{verbatim}
The default value used for $\alpha$, if initially $\alpha \le 0$, is
$\alpha = 0.5$. This means the median for default prior for the range,
is half of the length for which the PRW2 model is defined.

The 'step-size' $h$ is a typical distance between the locations where
process is defined and is $h=1$ in the definition above. If the
locations where the PRW2 process defined is in variable \texttt{loc},
then this parameter can be set by the user as
\begin{verbatim}
    h <- diff(range(loc)) / (length(loc) - 1)
\end{verbatim}
but does not need to! This is the value used for $h$, if initially
$h \le 0$.

The $\lambda$ parameter is computed from the calibration as defined
above if set to $\lambda \le 0$. If $\lambda > 0$, then this value of
$\lambda$ is used without computing the calibration, and the values of
$r_0$ and $\alpha$ are silently ignored.

Note that $r_0$ is in real scale, ie in the scale of \texttt{loc}, and
$r_0$ is \textbf{not} scaled by $h$.

Default values are \texttt{param = c(0, 0, 0, 0)}.

\subsection*{Example}

See the example in the documentation of PRW2 \texttt{inla.doc("prw2")}

\subsection*{Notes}

The details for this prior is available in the PhD thesis with link
HERE (ADD LINK LATER).



\end{document}


% LocalWords: 

%%% Local Variables: 
%%% TeX-master: t
%%% End: 
